% Generated draft: requires mathematical and source review.
\documentclass[11pt]{article}
\usepackage[margin=25mm]{geometry}
\usepackage{fontspec}
\setmainfont{Latin Modern Roman}
\usepackage{amsmath,amssymb,mathtools,unicode-math}
\setmathfont{Latin Modern Math}
\usepackage{enumitem,needspace,xurl,hyperref}
\setlength{\parindent}{0pt}
\setlength{\parskip}{5pt}
\setlength{\emergencystretch}{3em}
\newcommand{\N}{\mathbb{N}}
\newcommand{\Z}{\mathbb{Z}}
\newcommand{\Q}{\mathbb{Q}}
\newcommand{\R}{\mathbb{R}}
\newcommand{\C}{\mathbb{C}}
\newcommand{\F}{\mathbb{F}}
\newcommand{\A}{\mathbb{A}}
\newcommand{\PP}{\mathbb{P}}
\newcommand{\E}{\mathbb{E}}
\newcommand{\eps}{\varepsilon}
\newcommand{\dd}{\,\mathrm{d}}
\begin{document}
\section*{UnsolvedMath ID: 10000041}
\textbf{Source problem code:} catalogue entry
\subsection{Definitions and mathematical statement}
In the usual setting let \(G\) be an infinite connected locally finite vertex-transitive graph: its automorphism group acts transitively on vertices. A simple branching random walk has a branching genealogy whose children move independently by simple random-walk steps on \(G\), with a prescribed reproduction law. In a Galton--Watson formulation the offspring numbers are independent with a common law; that law is not specified in the excerpt.

Let \(\mathcal T\) be the trace graph, consisting of all visited vertices and traversed edges. The natural transience condition is that every fixed vertex is visited only finitely many times by the entire particle system, almost surely. This is stronger than transience of a single particle's path. An end of an infinite locally finite connected graph is an equivalence class of infinite rays, where two rays are equivalent if no finite vertex deletion separates their tails.

Under the intended trace interpretation, the conjecture is that on nonextinction a transient simple branching random walk has infinitely many ends:
\[
 \mathbb P\bigl(|\operatorname{Ends}(\mathcal T)|=\infty\mid\text{nonextinction}\bigr)=1
\]
for every admitted vertex-transitive graph and branching model. If the model guarantees survival, the conditioning is unnecessary. If extinction is possible, an unconditional assertion of infinitely many trace ends would omit finite extinct traces. The excerpt gives no branching-law or survival convention, so that issue must be resolved from the original.

It also does not explicitly identify the graph whose ends are meant. Trace ends, genealogy ends and ambient-graph ends are different objects; no conclusion about the latter two is substituted for the trace question. No Cayley, planarity or unimodularity restriction is inserted into ``any vertex-transitive graph''.

The source-side and bounded modern audits report no complete resolution, and the summary calls it open/triage. These are qualified source assessments rather than independently verified current status.
\subsection{Short English statement}
Should the trace of every transient branching random walk on a vertex-transitive graph have infinitely many ends on survival? The intended branching law, survival convention and trace interpretation need confirmation; ends of the genealogy or ambient graph are different questions.
\subsection{Sources}
\begin{itemize}
\item Local registry: \url{data/unsolvedmath/problems.json}, permanent UnsolvedMath ID \(10000041\), source code catalogue entry; supplied statement, background, and research summary, updated 2026-08-06.
\item Itai Benjamini, \emph{Coarse Geometry and Randomness}, Saint-Flour notes, Conjecture 9.33, PDF page 70: \url{https://www.wisdom.weizmann.ac.il/~itai/stflouraug24.pdf#page=70}. Preserved copy: \url{https://arquivo.pt/noFrame/replay/20201231041548id_/http://www.wisdom.weizmann.ac.il/~itai/stflouraug24.pdf}. Supplied extraction checked against rendered PDF, accessed 2026-07-29.
\item The provenance records cross-document deduplication, source-side status review and a bounded modern audit with no complete resolution located; no additional theorem reference is supplied.
\end{itemize}
\end{document}
